\hypertarget{mux1ee5c-tiuxeau-tuux1ea7n-1}{%
\subsubsection{Mục tiêu tuần 1:}}

\begin{itemize}
\tightlist
\item
  Làm quen với chương trình First Cloud AI Journey (FCAJ) và các thành
  viên trong nhóm.
\item
  Tìm hiểu các dịch vụ AWS cơ bản liên quan đến dự án.
\item
  Đọc hiểu kiến trúc tổng thể của dự án NewsRAG (phiên bản khóa luận gốc
  --- v1).
\item
  Clone repository và thiết lập môi trường phát triển local.
\end{itemize}

\hypertarget{cuxe1c-cuxf4ng-viux1ec7c-cux1ea7n-triux1ec3n-khai-trong-tuux1ea7n-nuxe0y}{%
\subsubsection{Các công việc cần triển khai trong tuần
này:}}

\begingroup
\small
\setlength{\tabcolsep}{3pt}
\renewcommand{\arraystretch}{1.15}
\begin{longtable}[]{@{}
  >{\raggedright\arraybackslash}p{(\columnwidth - 8\tabcolsep) * \real{0.0116}}
  >{\raggedright\arraybackslash}p{(\columnwidth - 8\tabcolsep) * \real{0.7248}}
  >{\raggedright\arraybackslash}p{(\columnwidth - 8\tabcolsep) * \real{0.0465}}
  >{\raggedright\arraybackslash}p{(\columnwidth - 8\tabcolsep) * \real{0.0581}}
  >{\raggedright\arraybackslash}p{(\columnwidth - 8\tabcolsep) * \real{0.1589}}@{}}
\toprule\noalign{}
\begin{minipage}[b]{\linewidth}\raggedright
Thứ
\end{minipage} & \begin{minipage}[b]{\linewidth}\raggedright
Công việc
\end{minipage} & \begin{minipage}[b]{\linewidth}\raggedright
Ngày bắt đầu
\end{minipage} & \begin{minipage}[b]{\linewidth}\raggedright
Ngày hoàn thành
\end{minipage} & \begin{minipage}[b]{\linewidth}\raggedright
Nguồn tài liệu
\end{minipage} \\
\midrule\noalign{}
\endhead
\bottomrule\noalign{}
\endlastfoot
2 & - Làm quen với các thành viên FCAJ - Đọc nội quy, quy định của
chương trình thực tập - Nhận nhiệm vụ và phân chia nhóm dự án &
02/06/2026 & 02/06/2026 & \\
3 & - Tìm hiểu tổng quan AWS và các nhóm dịch vụ: \qquad  + Compute
(EC2, ECS, Fargate, Lambda) \qquad  + Storage (S3) \qquad  + Database
(RDS, Aurora) \qquad  + Networking (VPC, SG) \qquad  + Messaging (SQS,
Kafka) & 03/06/2026 & 03/06/2026 &
\url{https://cloudjourney.awsstudygroup.com/} \\
4 & - Tạo AWS Free Tier account - Cài đặt AWS CLI \& cấu hình
credentials - Tìm hiểu AWS Management Console & 04/06/2026 & 04/06/2026
& \url{https://cloudjourney.awsstudygroup.com/} \\
5 & - Đọc tài liệu kiến trúc NewsRAG v1 (khóa luận gốc) - Tìm hiểu kiến
trúc RAG (Retrieval-Augmented Generation) - Nghiên cứu các thành phần:
Crawler, Kafka, ETL, Vector DB, LLM & 05/06/2026 & 05/06/2026 &
Pipeline\_v3.md \\
6 & - Clone repository NewsRAG về máy - Cài đặt Python virtual
environment - Chạy \texttt{docker-compose\ up\ -d} để khởi tạo
PostgreSQL + Kafka local - Đọc hiểu cấu trúc source code & 06/06/2026 &
06/06/2026 & README.md, Makefile \\
\end{longtable}
\endgroup

\hypertarget{kux1ebft-quux1ea3-ux111ux1ea1t-ux111ux1b0ux1ee3c-tuux1ea7n-1}{%
\subsubsection{Kết quả đạt được tuần
1:}}

\begin{itemize}
\tightlist
\item
  Hiểu được tổng quan về AWS và các nhóm dịch vụ liên quan đến dự án
  NewsRAG:

  \begin{itemize}
  \tightlist
  \item
    \textbf{Compute}: EC2, ECS Fargate (chạy container), Lambda
    (serverless functions)
  \item
    \textbf{Database}: RDS Aurora PostgreSQL (lưu trữ bài viết + vector
    search)
  \item
    \textbf{Messaging}: SQS (hàng đợi tin nhắn), Kafka (streaming ---
    dùng trong v1)
  \item
    \textbf{Container}: ECR (registry), ECS (orchestration)
  \end{itemize}
\item
  Đã tạo và cấu hình AWS Free Tier account thành công, bao gồm:

  \begin{itemize}
  \tightlist
  \item
    Access Key / Secret Key
  \item
    Region mặc định: \texttt{ap-southeast-2}
  \item
    AWS CLI hoạt động trên terminal
  \end{itemize}
\item
  Hiểu được kiến trúc tổng thể NewsRAG:

  \begin{itemize}
  \tightlist
  \item
    \textbf{v1 (Khóa luận)}: Lambda Crawler + Kafka + BAAI/bge-m3 +
    Qdrant Cloud + FastAPI/Next.js
  \item
    \textbf{v2 (Cải tiến)}: Fargate Crawler + SQS + Bedrock Titan Embed
    + Aurora pgvector
  \item
    Phân chia module: Crawl+Queue (A), Consumer+DB (B), ETL+Embed (C),
    RAG+Frontend (D)
  \end{itemize}
\item
  Clone repo thành công, thiết lập môi trường local:

  \begin{itemize}
  \tightlist
  \item
    Python 3.10 virtual environment
  \item
    PostgreSQL + Kafka chạy qua Docker Compose
  \item
    Đọc hiểu cấu trúc thư mục: \texttt{crawler/}, \texttt{consumer/},
    \texttt{etl/}, \texttt{vectorize/}, \texttt{search/},
    \texttt{database/}
  \end{itemize}
\end{itemize}
